\documentclass[12pt,letterpaper]{article}
\usepackage[utf8]{inputenc}
\usepackage[spanish]{babel}
\usepackage{amsmath}
\usepackage{amsfonts}
\usepackage{amssymb}
\usepackage[margin=2.5cm]{geometry}

\begin{document}

\begin{center}
    \Large \textbf{Evaluación de Examen 2: Unidad II} \\
    \large Física para Ciencias de la Salud (005-1713) \\
    \vspace{0.5cm}
    \textbf{Estudiante:} María Fernanda Malavé \quad \textbf{C.I.:} 32.585.180
    \vspace{0.3cm} \\
    \large \textbf{Calificación Final: 5,5/10 puntos}
\end{center}

\vspace{0.5cm}

\section*{Análisis de Resultados}

\subsection*{1. Cinemática (Aceleración media)}
\textbf{Procedimiento y resultado:} Extrae de manera estructurada los datos y plantea correctamente la ecuación de aceleración media. Sustituye y calcula $0,6 / 0,15$, arrojando el valor numérico exacto de $4$. Sin embargo, comete un error en el análisis dimensional al denotar el resultado con unidades de velocidad ($4 \text{ m/s}$) en lugar de aceleración ($4 \text{ m/s}^2$). \\
\textbf{Veredicto:} \textbf{Parcialmente correcto} (Penalización por unidades incorrectas) (1,5/2 pts).

\subsection*{2. Dinámica (Leyes de Newton)}
\textbf{Procedimiento y resultado:} Muestra un planteamiento inicial con la Segunda Ley de Newton, pero comete un error conceptual y algebraico muy grave: despeja la aceleración afirmando que es el producto de la masa por la fuerza ($a = m \cdot F$) en lugar del cociente ($a = \frac{F}{m}$). Este fallo procedimental la lleva a calcular $25 \cdot 150 = 3750$ y a indicarlo erróneamente en $\text{m/s}$. \\
\textbf{Veredicto:} \textbf{Incorrecto} (Se otorgan puntos parciales por extracción de datos) (0,5/2 pts).

\subsection*{3. Trabajo Mecánico}
\textbf{Procedimiento y resultado:} Extrae los datos y plantea la fórmula del trabajo. Sin embargo, al aplicar el factor trigonométrico asume un ángulo incorrecto y arbitrario, calculando un factor de $0,76$ en lugar del $\cos(0^\circ) = 1$ correspondiente a un desplazamiento paralelo a la fuerza. Esto la conduce al resultado errado de $243,2 \text{ J}$ en lugar de los $320 \text{ J}$ reales. \\
\textbf{Veredicto:} \textbf{Incorrecto} (Se otorgan puntos parciales por extracción de datos y ecuación base) (0,5/2 pts).

\subsection*{4. Energía Potencial}
\textbf{Procedimiento y resultado:} Extrae rigurosamente los datos del problema y plantea a la perfección la ecuación de la energía potencial gravitatoria ($E_p = m \cdot g \cdot h$). Sustituye correctamente los factores en el papel ($1,2 \text{ kg} \cdot 9,8 \text{ m/s}^2 \cdot 1,5 \text{ m}$), pero comete un error aritmético drástico al efectuar la multiplicación, obteniendo $58,8 \text{ J}$ en lugar del valor exacto de $17,64 \text{ J}$ (matemáticamente parece haber multiplicado por $5$ en lugar de $1,5$). \\
\textbf{Veredicto:} \textbf{Parcialmente correcto} (Desarrollo teórico excelente pero resultado aritmético equivocado) (1/2 pts).

\subsection*{5. Potencia Mecánica}
\textbf{Procedimiento y resultado:} Extrae los datos y relaciona correctamente el trabajo efectuado y el tiempo planteando la fracción respectiva ($P = \frac{W}{t}$). El desarrollo de la división de $75 \text{ J} / 60 \text{ s}$ arroja la cifra correcta y verificada de $1,25 \text{ W}$, con un buen manejo de unidades. \\
\textbf{Veredicto:} \textbf{Correcto} (2/2 pts).

\vspace{0.5cm}
\section*{Comentarios Generales y Sugerencias}
La estudiante demuestra poseer una muy buena organización y claridad metodológica, separando siempre la zona de \textit{Datos} y el desarrollo con recuadros, lo cual es muy positivo. No obstante, es imperativo reforzar el aspecto teórico y aritmético.
\begin{itemize}
    \item Se requiere un profundo repaso sobre los **despejes algebraicos** de fórmulas base (notable en la Pregunta 2). Multiplicar variables en lugar de dividirlas distorsiona por completo la realidad del fenómeno físico.
    \item Es importante revisar dos veces los cálculos efectuados en la calculadora o a mano para evitar errores de tipeo que arruinen planteamientos perfectos (como ocurrió en la Pregunta 4).
    \item Vigilar cuidadosamente el **análisis dimensional**: la aceleración se mide en $\text{m/s}^2$. Las unidades son fundamentales en las ciencias biomédicas para evitar diagnósticos equívocos.
\end{itemize}

\end{document}